\documentclass[../../main.tex]{subfiles}
% \documentclass[a4paper,10pt]{ltjsarticle}
\begin{document}
\rhead{京都大学後期$2003$年$2$番}
\setcounter{equation}{0}
\begin{oframed}
  \documentclass[../../main.tex]{subfiles}
\begin{document}

一辺の長さが1の正三角形$\mathrm{ABC}$の辺$\mathrm{AC}$上に点$\mathrm{D}$をとり，線分$\mathrm{BD}$に沿ってこの三角形を折り曲げ，4点$\mathrm{A}$，$\mathrm{B}$，$\mathrm{C}$，$\mathrm{D}$を頂点とする四面体を作り，その体積を最大にすることを考える．
体積が最大になる時の$\mathrm{D}$の位置と，その時の四面体の体積を求めよ．

\end{document}

\end{oframed}
\setlength{\columnseprule}{0.4pt}
\begin{multicols}{2}
  {\bf [解答]}
  三角形ABCにおいて
  \begin{align}
    & \angle ADB = \theta & \left( \dfrac{\pi}{3} < \theta < \dfrac{2\pi}{3} \right) \label{eq:1}
  \end{align}
  とおく．
  C から直線 BD に下ろした垂線 H とおく．

  \begin{figure}[H]
    \centering
    \begin{tikzpicture}
      % Define vertices of the equilateral triangle ABC
      \coordinate (A) at (0, 3.464); % Height of equilateral triangle with side length 4
      \coordinate (B) at (-2, 0);
      \coordinate (C) at (2, 0);

      % Draw the equilateral triangle ABC
      \draw (A) -- (B) -- (C) -- cycle;

      % Mark angles of equilateral triangle (optional, just for visual confirmation)
      \pic [draw, angle eccentricity=1.5, angle radius=0.7cm, "$\frac{\pi}{3}$"] {angle = B--A--C};
      % \pic [draw, angle eccentricity=1.2, angle radius=0.7cm, "$\frac{\pi}{3}$"] {angle = A--B--C};
      % \pic [draw, angle eccentricity=1.2, angle radius=0.7cm, "$\frac{\pi}{3}$"] {angle = B--C--A};

      % Point D on AC such that AD:DC = 3:1
      % Since AC has length 4, AD = 3 and DC = 1
      \coordinate (D) at ($(A)!0.6!(C)$); % 0.75 along the path from A to C

      % Draw line BD
      \draw (B) -- (D);

      % Draw the perpendicular from C to BD, let H be the foot of the perpendicular
      \path (C) -- ($(B)!(C)!(D)$) coordinate (H);
      \draw[dashed] (C) -- (H) -- (D);

      % Mark the right angle at H
      % \draw (H) rectangle ($(H) + (-0.2, 0.2)$); % Adjust coordinates if needed for better placement

      % Labels
      \node[above] at (A) {A};
      \node[left] at (B) {B};
      \node[right] at (C) {C};
      \node[below] at (D) {D};
      \node[above] at (H) {H};

      % Mark angle ADB as theta
      \pic [draw, angle eccentricity=2.0, angle radius=0.5cm, "$\theta-\frac{\pi}{3}$"] {angle = C--B--D}; % Angle DBC is theta
      \pic [draw, angle eccentricity=1.5, angle radius=0.5cm, "$\theta$"] {angle = A--D--B}; % Angle DBC is theta

      \fill (D) circle (0.05);
      \fill (H) circle (0.05);
      % Add other lines from the reference image, if desired (e.g., lines from A to H)
      % \draw[dashed] (A) -- (H);
    \end{tikzpicture}
    \caption{三角形ABCと点D, Hの様子}
    \label{2003-2:fig:1}
  \end{figure}

  \begin{figure}[H]
    \centering
    \begin{tikzpicture}[x={(1cm,0cm)}, y={(0.55cm,0.75cm)}, z={(0cm,1cm)},
      scale=0.95, >=stealth, every node/.style={font=\small}]
      % 一辺4の正三角形 ABC で D が AC を 3:7 に内分する例．直線 BD を横軸，面 ABD を床面 z=0 とする．
      % 面 BCD を BD で折って面 ABD に垂直に立てると，C は H の真上に来る．
      \coordinate (B)  at (0,0,0);
      \coordinate (D)  at (3.5553,0,0);
      \coordinate (A)  at (3.8253,1.1692,0);
      \coordinate (H)  at (2.9252,0,0);        % C から直線 BD への垂線の足
      \coordinate (C0) at (2.9252,-2.7282,0);  % 折る前の C
      \coordinate (C)  at (2.9252,0,2.7282);   % 折った後の C（|CH| は不変）
      % 床面（三角形 ABD）
      \fill[pattern=north east lines, pattern color=blue!45!black] (A) -- (B) -- (D) -- cycle;
      % 折る前の三角形 BCD（参考）
      \draw[densely dotted, thick, gray!70!black] (B) -- (C0) -- (D);
      \draw[densely dotted, thick, gray!70!black] (C0) -- (H);
      \draw[->, gray!70!black] ($(C0)+(-0.35,0,0)$) to[bend left=45]
      node[pos=0.2, left, font=\footnotesize] {折る} ($(C)!0.55!(H)+(-0.3,0,0)$);
      % 垂線 CH と直角の印（CH ⟂ BD，CH ⟂ 床面）
      \draw[thick, red!85!black] (C) -- (H);
      \draw[red!85!black] (H) ++(-0.25,0,0) -- ++(0,0,0.25) -- ++(0.25,0,0);
      % 四面体 ABCD（辺 AB は手前の面 BCD に隠れる）
      \draw[line width=1.3pt, blue!85!black] (B) -- (D) -- (A);
      \draw[line width=1.3pt, blue!85!black, densely dashed] (A) -- (B);
      \draw[line width=1.3pt, teal!85!black] (C) -- (B) (C) -- (D) (C) -- (A);
      \foreach \p in {A, B, D, H, C, C0} {
        \node[circle, fill=black, inner sep=1.2pt] at (\p) {};
      }
      \node[right=1pt] at (A) {$\mathrm{A}$};
      \node[left=1pt] at (B) {$\mathrm{B}$};
      \node[below right=1pt] at (D) {$\mathrm{D}$};
      \node[above right=1pt] at (H) {$\mathrm{H}$};
      \node[above=1pt] at (C) {$\mathrm{C}$};
      \node[below=1pt, gray!70!black] at (C0) {$\mathrm{C}$（折る前）};
    \end{tikzpicture}
    \caption{面BCDをBDで折り，面ABDに垂直に立てた時の四面体ABCD（$CH$は面ABDに垂直）}
    \label{2003-2:fig:2}
  \end{figure}

  四面体 $ABCD$ の体積が最大の時
  の $\triangle ABD$ を底面とみた時の高さは $CH$ である．

  従って，四面体ABCDは底面が三角形ABD，高さCHの錐体で，
  その体積$V(\theta)$は，
  \begin{align}
    V(\theta)
    & = \dfrac{1}{3}(\triangle ABD) |CH|          \\
    & = \dfrac{1}{6}|CH||AD||AB|\sin\dfrac{\pi}{3} \\
    & = \dfrac{\sqrt{3}}{12}|CH||AD||AB|
  \end{align}
  である．

  以下CH,AD,ABの長さを計算する．
  まず，題意より三角形ABCは一辺の長さ$1$だから
  \begin{align}
    |AB| = 1
  \end{align}
  である．

  次にADは，三角形ABDに正弦定理を用いて
  \begin{align}
    & \dfrac{|AD|}{\sin \angle ABD } = \dfrac{|AB|}{\sin \angle ADB} \\
    & |AD| = \dfrac{\sin\left(\dfrac{2\pi}{3}-\theta\right)}{\sin\theta}
  \end{align}
  である．

  最後にCH については，三角形BCHを考えて
  \begin{align}
    |CH|
    & = |BC|\sin \angle HBC  \\
    & = \sin\left(\theta - \dfrac{\pi}{3}\right)
  \end{align}
  である．

  以上から，四面体 $ABCD$ の体積 $V(\theta)$ は
  \begin{align}
    V(\theta)
    & = \dfrac{\sqrt{3}}{12}|CH||AD||AB|                                                      \\
    & = \dfrac{\sqrt{3}}{12} \dfrac{\sin\left(\dfrac{2\pi}{3}-\theta\right)}{\sin\theta} \sin(\theta - \dfrac{\pi}{3})
  \end{align}
  である．
  積和公式を用いてさらに変形して
  \begin{align}
    V(\theta)
    & = \dfrac{\sqrt{3}}{12} \dfrac{-1}{2}\dfrac{\cos\dfrac{\pi}{3}-\cos(\pi-2\theta)}{\sin\theta} \\
    & = \dfrac{-\sqrt{3}}{24} \dfrac{\dfrac{1}{2}+\cos 2\theta}{\sin\theta}
  \end{align}
  さらに倍角公式から$\cos 2\theta = 1-2\sin^2\theta$だから，全体を$\sin\theta$だけでかけて
  \begin{align}
    V(\theta)
    & = -\dfrac{\sqrt{3}}{48} \dfrac{1+2(1-2\sin^2\theta)}{\sin\theta}        \\
    & = -\dfrac{\sqrt{3}}{48} \dfrac{3-4\sin^2\theta}{\sin\theta}             \\
    & = \dfrac{\sqrt{3}}{48} \left(4\sin\theta - \dfrac{3}{\sin\theta}\right)
  \end{align}
  となる．
  ここで，\cref{eq:1}から，区間内で$\dfrac{\sqrt{3}}{2}<\sin\theta\le 1$であり，
  この区間で$\sin\theta$も$-1/\sin\theta$も$\sin\theta$に関する単調増加関数だから，
  $V(\theta)$は$\sin\theta$ の単調増加関数で，$\sin\theta=1$ すなわち $\theta=\pi/2$
  の時に最大値を取る．
  この時$D$ が $AC$ の中点の時となる．
  求める最大値は
  \begin{align}
    \max V
    & = V\left(\dfrac{\pi}{2}\right)
    = \dfrac{\sqrt{3}}{48}
  \end{align}
  である．$\cdots$(答)

  \vspace{10pt}
  {\bf [解説]}

  \newpage
\end{multicols}
\end{document}
